\section{Coupling the finite and independent tests}
\label{sec:coupling}

Fix a graph $G$ satisfying Proposition~\ref{prop:prefix}.
The independent probabilities from Proposition~\ref{prop:unfolding}
approximate survival only if repeated uses of a finite triangle priority
are sufficiently rare.  We first give an exact coupling and then reduce
its failure to a graph configuration supported on two visited call paths.
Throughout this section, a \emph{triangle type} means a member of $\cT$,
as distinguished from its possibly many occurrences in an unfolding.

Fix $k\in\{1,2\}$.  Sample $k$ independent uniform oriented edges of $G$,
with replacement and independently of all priorities.  There are $2m$
possible oriented edges.  The orientation names the two endpoints and
does not affect the test.  With these roots, compare two experiments:
the finite tests sharing one priority for each triangle type, and the
independent unfoldings with disjoint occurrence trees for different
roots.  Every root has threshold $1$.  Execute the roots in a fixed
order, including every root regardless of the preceding answers.

At each visited call, expose the priorities of \emph{all} candidates in
a fixed type order before sorting them for the test.  Thus an exposed
list includes candidates above the threshold and candidates never
traversed because an earlier child returned true.  In the independent
experiment, let $\mathcal C_k$ be the event that a triangle type appears
at least twice among all these exposed lists.  There is no repetition
within one list: a candidate meeting both edges of a pair focus would
be its omitted parent triangle.

\begin{lemma}[Coupling by exposed triangle types]
\label{lem:coupling}
Let $X^{\mathrm{fin}},X^{\mathrm{ind}}\in\{0,1\}^k$ be the vectors of
root answers in the two experiments, where $1$ denotes true.
For every $H\subseteq\{0,1\}^k$,
\[
 \bigl|\Prb(X^{\mathrm{fin}}\in H)
       -\Prb(X^{\mathrm{ind}}\in H)\bigr|
 \le \Prb_{\mathrm{ind}}(\mathcal C_k).
\]
\end{lemma}

\begin{proof}
Assign independent uniform priorities $U_T$, $T\in\cT$, for the finite
experiment, and an independent supply of fresh uniform variables.
Run the unfolding tests as follows.  At the first exposure of a type
$T$, use $U_T$; at each later exposure of that type, use a fresh
variable instead.

This procedure has the independent-unfolding law.  Indeed, the type
queried next is determined by $G$, the root labels and the execution
history.  If $Q$ is the set of types exposed so far, that history
depends only on $(U_T:T\in Q)$ and on the fresh variables already
used.  Every recursive computation contributing to a returned answer
is part of this history, with its candidate lists exposed explicitly.
Conditional on the history, the variables $(U_T:T\notin Q)$ are still
independent uniforms.  A fresh variable used for a repeated type is
independent of the same history.  Induction over exposures therefore
gives independent uniform occurrence priorities.  The unused
occurrence variables may be filled in independently afterward.

The evaluations terminate almost surely by Section~\ref{sec:queries}.
Until a repeated type is exposed, the unfolding and finite executions
use the same priority at every candidate.  They consequently have the
same sorted lists, recursive calls, thresholds, answers and stopping
decisions.  On $\mathcal C_k^c$ this agreement holds for the whole
execution of all roots.  The stated inequality follows from this
coupling of the answer vectors.
\end{proof}

\subsection{The graph of visited calls}

We now work entirely in the independent experiment.  Its visited calls
form a finite rooted forest almost surely.  Construct a graph $\Gamma$
and a label map $\lambda:V(\Gamma)\to V(G)$ from this forest.
Each root contributes two new \emph{formal} vertices $x,y$, the edge
$xy$, and labels given by its sampled ordered endpoints.  Distinct
roots have disjoint formal vertices even if their labels coincide.
If a visited child is attached through a formal focus edge $xy$ of its
parent, introduce a new formal vertex $z$, label it by the third
vertex of the child's triangle type, and add the edges $xz,yz$.
These two new edges are the child's focus.  If $xy$ is ordered, order
them as $(x,z),(y,z)$, so subsequent choices of focus edge have fixed
names.

The \emph{birth call} of a formal vertex or edge is the call that
introduced it.  Order the construction with all roots first and every
parent before its children.  The following properties hold in any such
order:
\begin{enumerate}[label=(\roman*)]
\item Every formal edge has adjacent labels in $G$, and the three
vertices in any one gluing have distinct labels.
\item Every formal edge is focused at exactly one call, its birth
call.  An attachment edge is not in the child's focus.
\item No step adds an edge between two vertices already present.
\item The only triangles of $\Gamma$ are the triangles introduced by
the gluings.
\end{enumerate}
For the last property, consider the latest vertex of a formal
triangle.  The initial disjoint root edges contain no triangle.  A
non-root vertex has exactly two older neighbours, namely the endpoints
of its attachment edge, so the triangle must be its gluing triangle.

Figure~\ref{fig:call-graph} shows the distinction between a visited call
and an exposed candidate. The formal graph records the former, but an
extra adjacency between its labels can account for repetitions among
the latter.

\begin{figure}[htbp]
\centering
\begin{tikzpicture}[x=1cm,y=1cm,
 call/.style={draw=black!65,rounded corners=2pt,fill=white,
               inner sep=5pt,font=\small},
 vertex/.style={circle,fill=black,inner sep=1.6pt},
 every label/.style={font=\small}]
\node[font=\small\bfseries] at (0,.6) {Visited call path};
\node[call] (root) at (0,0) {focus $\{ab\}$};
\node[call] (first) at (0,-1.25) {focus $\{ac,bc\}$};
\node[call] (second) at (0,-2.5) {focus $\{ad,cd\}$};
\draw[-{Stealth},black!65] (root)--node[left,font=\small] {$abc$} (first);
\draw[-{Stealth},black!65] (first)--node[left,font=\small] {$acd$} (second);
\node[font=\small,align=left,anchor=west] at (1.5,0)
 {$abd$ exposed\\through $ab$};
\node[font=\small,align=left,anchor=west] at (1.5,-2.5)
 {$abd$ exposed\\through $ad$};
\begin{scope}[xshift=6.5cm]
\node[font=\small\bfseries] at (0,.6) {Formal graph};
\coordinate (a) at (-1.35,-1.25);
\coordinate (b) at (0,0);
\coordinate (c) at (1.35,-1.25);
\coordinate (d) at (0,-2.5);
\fill[black!3] (a)--(b)--(c)--cycle;
\fill[linkblue!7] (a)--(c)--(d)--cycle;
\draw[black!65] (a)--(b)--(c)--(a)--(d)--(c);
\draw[dashed,linkblue,thick] (b)--(d);
\node[vertex,label=left:$a$] at (a) {};
\node[vertex,label=above:$b$] at (b) {};
\node[vertex,label=right:$c$] at (c) {};
\node[vertex,label=below:$d$] at (d) {};
\end{scope}
\end{tikzpicture}
\caption{A portion of the visited construction, with distinct labels
written as the formal vertex names. The left arrows are parent--child steps:
traversing $abc$ introduces $c$, and traversing $acd$ introduces $d$.
The solid edges on the right are the corresponding edges of $\Gamma$.
The dashed segment $bd$ denotes an edge between labels in $G$ that is
absent from $\Gamma$. The type $abd$ is exposed at the root through $ab$
and at the last displayed call through $ad$. The collision is therefore detected from
these two candidate lists, regardless of whether either occurrence is
traversed. Other calls are not shown.}
\label{fig:call-graph}
\end{figure}

For a formal vertex $a$, its \emph{ancestral call path} is the path
from a root to the birth call of $a$, including both ends.  Given two
formal vertices $a,b$, let $\Gamma_{a,b}$ be the graph constructed from
the union of their ancestral call paths, using the root edges and
gluings on those paths.  All its calls are visited.  We call $(a,b)$
an \emph{extra-edge witness} if $a,b$ are distinct and nonadjacent in
$\Gamma$, their labels are adjacent in $G$, and $\lambda$ is injective
on $V(\Gamma_{a,b})$.

\begin{lemma}[A geometric witness for a collision]
\label{lem:witness}
If the initial root labels are all distinct, then every occurrence of
$\mathcal C_k$ produces an extra-edge witness.  Moreover, uniformly
over the graphs under consideration,
\begin{equation}
\label{eq:collision-witness}
 \Prb_{\mathrm{ind}}(\mathcal C_k)
 \le O(n^{-1})+
 \Prb_{\mathrm{ind}}(\text{an extra-edge witness exists}).
\end{equation}
The $O(n^{-1})$ term may be omitted when $k=1$.
\end{lemma}

\begin{proof}
First suppose that $\lambda$ is injective on all of $\Gamma$.
If an edge of $G$ between two labels is absent from the formal graph,
its two formal endpoints give the required witness.  It remains to
show that a collision is impossible when there is no such extra edge.

Suppose that a triangle type $T$ is exposed at two calls.  It cannot
occur twice in one list, so the exposing calls are distinct.  Their
responsible focus edges are distinct by property~(ii), and their
images are distinct by injectivity.  Those two edges already represent
all three vertices of $T$, even if neither candidate occurrence is
traversed.  Since there are no extra edges between labels, the three
formal vertices form a triangle of $\Gamma$.  By property~(iv), it is
a gluing triangle.  Its two new edges are focused only at that
gluing's child call, which omits the triangle as its parent.  Its
attachment edge is focused at only one call.  Hence the triangle can
appear in at most one candidate list, a contradiction.

Next suppose that the labels are not injective.  Choose the first
repetition in a construction order with roots first and parents
before children.  The distinct-root assumption means that this
repetition occurs when a new vertex $z$ is attached through an edge
$xy$.  Write $w$ for the earlier vertex with $\lambda(w)=\lambda(z)$.
All earlier labels are distinct.  Property~(i) shows that $x,y,w$ are
distinct and that $\lambda(x)\lambda(w)$ and
$\lambda(y)\lambda(w)$ are edges of $G$.

If either $xw$ or $yw$ is absent formally, it gives an extra edge
among the earlier vertices.  The ancestral paths of its endpoints
are entirely earlier in the construction, so their labels are
injective.  Property~(iii) ensures that the pair remains nonadjacent
in the completed graph.  This is the desired witness.

Otherwise $xyw$ is a formal triangle and hence an earlier gluing
triangle.  There are two possibilities for its latest vertex.
If it is $x$ or $y$, then $xy$ was born at this gluing.  The current
call focusing $xy$ is exactly its birth call, and omits the triangle
$xyw$ as its parent.  It therefore cannot have the proposed child
with the same triangle type.  If the latest vertex is $w$, then $w$
was attached through $xy$.  The children introducing $w$ and $z$
would belong to the same unique call focusing $xy$, and would have
the same triangle type.  That call contains this candidate only once
and traverses it at most once, again a contradiction.  Thus a first
repeated label always produces an earlier extra edge.

Finally, either endpoint of a uniform oriented root has distribution
\[
 \pi(v)=\frac{\deg_G(v)}{2m},
 \qquad \max_v\pi(v)=O(n^{-1}),
\]
by the degree estimate and $2m=n^2p$.  For two independent roots, any
fixed pair of their endpoints coincides with probability
$\sum_v\pi(v)^2\le\max_v\pi(v)$.  There are four pairs to consider;
the two endpoints within each root are distinct.  A union bound
therefore gives the term $O(n^{-1})$ in
\eqref{eq:collision-witness}.
\end{proof}

The root experiment still samples with replacement; we have only
charged its overlaps to the error bound.  The remaining task is to
bound the probability of a visited, injectively labelled pair of call
paths whose two marked vertices have an extra edge.  This uses both
the template estimates for $G$ and the independent-test probabilities.
